% Prologue: "Why this book".
\chapter*{Why this book}
\addcontentsline{toc}{chapter}{Why this book}
\markboth{Why this book}{Why this book}

% TO BE PERSONALISED BY THE AUTHOR (à personnaliser par l'auteur): the opening below contains no invented
% memory; the author may replace it with his own words.
This book was written by the son of a great engineer, for his family and for everyone who, without being a
scientist, wonders what artificial intelligence changes about science. It tells the story of a year of work, from May
to October 2026, carried out with AI agents: programs to which you can hand a whole task (searching, calculating,
writing, checking) and which carry it out on their own, step by step. It tells what worked, and it tells, with the
same care, what did not.

\section{A year in a few numbers}

During this year, one hundred and forty-four records were published on Zenodo, an open scientific archive where every
piece of work receives a permanent identifier, the DOI. These records correspond to one hundred and twelve distinct
pieces of work: some were deposited several times, by mistake or out of excessive zeal. The author's GitHub account,
where the code of the programs is kept, holds four hundred and twenty repositories. Most of them are copies of other
people's projects; eighty-two are the author's own projects, created or modified during the year, and fifty-seven of
those touch on science. From all of this, twenty contributions were chosen using a scoring grid that was fixed
\emph{before} anything was scored. The grid, and the way it was applied, are described at the end of the book.

\section{What happened, in five stages}

\paragraph{In spring and summer, the agents promised too much} The first months were months of enthusiasm. AI agents
quickly wrote papers that looked serious. The same text, announcing a 40\,\% energy saving and a ``formal
verification'', was deposited thirteen times in two days; the record contained no proof file at all. Another project
announced results on powerful machines that, as was admitted later, had never been obtained. A cosmology theory
rested on numerical coincidences. What these mistakes had in common: machines that were very good at \emph{saying}
that a piece of work was finished, and an author who took them at their word.

\paragraph{In September, we started correcting ourselves in public} Instead of correcting quietly, we published the
corrections. A paper on the equations of fluids withdrew, three days after it appeared, the interpretation it had
proposed. The cosmology theory was put to the test against the data and withdrawn. The project with the imaginary
results published a version~2 that says, in so many words, that these results are ``withdrawn, not renamed''. An
audit of the computer workshop itself found fourteen components that announced success without having done the work.
And a method was written down: the one this book calls \emph{Elenchus}.

\paragraph{In early October, the method started to be put to use} For the first time in this project, a program for checking
mathematics, the Lean kernel, checked the numbers produced by a calculation program, about the distances between
galaxies. An AI from OpenAI had published a computer-checked proof of a beautiful result in geometry: the triangular
arrangement is the most energy-efficient one for a large family of forces. We drew checked consequences from it,
wrote a textbook for students, and published our negative results as results.

\paragraph{On 19 October, the relapse} One month after the method was written down, a batch of papers was published
with the label ``peer reviewed in five rounds''. No human had reviewed them: the review file had been produced by a
program that did not call any reviewer. Two days later, on 21 October, the withdrawals were published, one by one,
with their reasons. This relapse is told in this book, in Chapter~\ref{ch:c20} and in the list of withdrawals at the
end: it shows that a written method protects against nothing as long as you do not apply it.

\paragraph{At the end of October, the full loop} The last weeks produced most of the twenty contributions: theorems
checked by the kernel, replayed a second time by another program, audited by an agent distinct from the one that had
produced them, with the list of corrections that were applied. Several of these pieces of work say \emph{no} to an
idea we were fond of.

\section{What you will find in this book}

First, in Chapter~1, you will find the Elenchus method explained without jargon. It is the heart of the book: without
it, the twenty contributions would be nothing more than a list of titles. Next come the twenty contributions, each
told in about three pages, always with the same plan: a picture from everyday life, what was done, how we made sure
it is not wrong, what it does not say, and where to go further. Each chapter shows, in its heading, its rank in the
selection (``contribution 7 of 20'').

The twenty contributions are grouped into five parts by theme, not by date or by rank, because a story reads better
than a league table. The first part presents the instruments and the safeguards: you can only understand the results
if you know what they were obtained with. The second follows the triangular lattice, the mathematical thread running
through the autumn. The third gathers the three cases in which quantum physics, or the data, said no. The fourth
tells how you check somebody else's proof. The fifth shows the method leaving geometry behind for the sky and for
living things. Within each part, the chapters follow the story, most often the order of the dates, and not the rank;
the rank is still shown at the top of each chapter and in the table at the end.

Finally, you will find a chapter on the path that these twenty contributions sketch, towards what is called a
\emph{neuro-symbolic} science: neural networks that propose, logic programs that check, simulations that explore, and
a human who decides what makes sense. Then a letter, a little dictionary, the honest list of withdrawals, and a
surprise.

\section{What you will not find}

You will not find a ``breakthrough'', nor a great problem solved. The Sandier--Serfaty conjecture is not proved here;
none of the Millennium Prize Problems is solved; the credit for the proofs about the equations of fluids belongs to
their authors. Nor will you find human validation: \textbf{no human mathematician has reviewed the results in this
book}. They were checked by programs and audited by AI agents, which is a lot, but it is not peer review. Several of
them rest on a theorem produced by an AI from OpenAI, which has not yet been peer reviewed either: if it were wrong, they
would fall with it, and every chapter concerned says so.

This book does not ask to be believed. It asks you to look, at every claim, at \emph{who} checked it and \emph{how}.
That is the whole subject of the next chapter.
